\documentclass[11pt]{article}
\usepackage[margin=1in]{geometry}
\usepackage{amsmath,amssymb}
\newcommand{\finalanswer}[1]{\par\medskip\noindent\fbox{\parbox{0.94\linewidth}{\textbf{Answer:} #1}}}
\title{STAT 412 Homework 6 --- Question 12}
\author{}\date{}
\begin{document}\maketitle
\section*{Problem}
In a survey of $2029$ adults, $717$ made a New Year's resolution to eat healthier. Construct 90\% and 95\% confidence intervals for the population proportion.
\section*{Work}
\[
\hat p=\frac{717}{2029}\approx0.3534.
\]
For a proportion,
\[
\hat p\pm z_{\alpha/2}\sqrt{\frac{\hat p(1-\hat p)}{n}}.
\]
For 90\%, $z=1.645$:
\[
E\approx1.645\sqrt{\frac{0.3534(0.6466)}{2029}}\approx0.0175,
\]
so the interval is
\[
(0.3359,0.3708)\approx(0.336,0.371).
\]
For 95\%, $z=1.960$:
\[
E\approx0.0208,\qquad (0.3326,0.3742)\approx(0.333,0.374).
\]
The interpretation uses the population proportion, not the sample proportion: with the given confidence, the population proportion is between the endpoints of the interval. The 95\% interval is wider because it uses a larger critical value than the 90\% interval.
\finalanswer{90\% CI $(0.336,0.371)$; 95\% CI $(0.333,0.374)$. Fill-ins: population proportion; between the endpoints. The 95\% interval is wider.}
\end{document}
